% LR 调度策略对比：横轴 epoch，纵轴 LR
% 4 条曲线：常数 / 指数衰减 / Cosine annealing / Cosine with warmup
\documentclass[tikz,border=4pt]{standalone}
\usepackage{ctex}
\usepackage{amsmath}
\usetikzlibrary{calc, arrows.meta}

\begin{document}
\begin{tikzpicture}[scale=1.0, thick]

% 参数定义
% 横轴：epoch 0..100，映射到 x = 0..8
% 纵轴：LR 0..1（相对），映射到 y = 0..4

% 坐标轴
\draw[->, gray] (-0.2, 0) -- (9.0, 0) node[right, font=\small] {Epoch};
\draw[->, gray] (0, -0.2) -- (0, 4.5) node[above, font=\small] {学习率 $\eta_t$};

% 刻度：横轴
\foreach \x/\lab in {0/0, 2/20, 4/40, 6/60, 8/80, 8.5/100}{
  \draw[gray, thin] (\x, 0) -- (\x, -0.07);
  \node[below, font=\footnotesize, gray] at (\x, -0.1) {\lab};
}
% 纵轴刻度
\foreach \y/\lab in {0/0, 1/0.25, 2/0.5, 3/0.75, 4/1.0}{
  \draw[gray, thin] (0, \y) -- (-0.07, \y);
  \node[left, font=\footnotesize, gray] at (-0.1, \y) {\lab};
}

% ── 曲线1：常数 LR（黑色虚线）───────────────────────────
\draw[black!60, dashed, thick] (0, 2.4) -- (8.5, 2.4);
\node[right, font=\footnotesize, black!60] at (8.5, 2.4) {};

% ── 曲线2：指数衰减（绿色）──────────────────────────────
% eta_t = eta_0 * 0.95^t, eta_0=0.6, t=[0..100]
% x=t/100*8.5, y=0.6*0.95^t*4 (scaled)
\draw[green!60!black, very thick]
  plot[smooth, domain=0:8.5, samples=80]
  (\x, {4 * 0.6 * exp(-\x / 8.5 * 100 * 0.05126)});
% 0.05126 = -ln(0.95), so eta_t = 0.6 * exp(-0.05126*t)

% ── 曲线3：Cosine annealing（蓝色）─────────────────────────
% eta_t = 0.5 * eta_max * (1 + cos(pi*t/T))
% eta_max=0.8, T=100, x=t/100*8.5, y=0.5*0.8*(1+cos(pi*t/100))*4
\draw[blue, very thick]
  plot[smooth, domain=0:8.5, samples=200]
  (\x, {4 * 0.5 * 0.8 * (1 + cos(3.14159 * \x / 8.5))});

% ── 曲线4：Cosine with warmup（红色）────────────────────────
% warmup: epoch 0..10 (x=0..0.85), linear 0 to 0.8
% cosine: epoch 10..100 (x=0.85..8.5)
\draw[red, very thick]
  plot[smooth, domain=0:0.85, samples=30]
  (\x, {4 * 0.8 * \x / 0.85});
% Warmup 阶段终点
\filldraw[red] (0.85, 3.2) circle (2pt);
% Cosine 阶段
\draw[red, very thick]
  plot[smooth, domain=0.85:8.5, samples=200]
  (\x, {4 * 0.5 * 0.8 * (1 + cos(3.14159 * (\x - 0.85) / 7.65))});

% Warmup 区域标注
\draw[red!50, thin, dashed] (0.85, 0) -- (0.85, 3.5);
\node[above, font=\footnotesize, red!70] at (0.43, 3.35) {Warmup};
\draw[<->, red!50, thin] (0, 3.5) -- (0.85, 3.5) node[midway, above, font=\tiny, red!50] {10e};

% ── 图例 ──────────────────────────────────────────────────
\begin{scope}[xshift=1.0cm, yshift=0.2cm]
  % 常数
  \draw[black!60, dashed, thick] (0, 3.8) -- (1.0, 3.8);
  \node[right, font=\footnotesize, black!60] at (1.0, 3.8) {常数 LR};

  % 指数衰减
  \draw[green!60!black, very thick] (0, 3.4) -- (1.0, 3.4);
  \node[right, font=\footnotesize, green!60!black] at (1.0, 3.4) {指数衰减};

  % Cosine
  \draw[blue, very thick] (0, 3.0) -- (1.0, 3.0);
  \node[right, font=\footnotesize, blue] at (1.0, 3.0) {Cosine Annealing};

  % Cosine + warmup
  \draw[red, very thick] (0, 2.6) -- (1.0, 2.6);
  \node[right, font=\footnotesize, red] at (1.0, 2.6) {Cosine + Warmup};

  % 图例边框
  \draw[black!30, thin, rounded corners=2pt] (-0.2, 2.4) rectangle (3.6, 4.0);
\end{scope}

% ── 标题 ──────────────────────────────────────────────────
\node[font=\small\bfseries] at (4.25, -0.6) {学习率调度策略对比（归一化幅度）};

\end{tikzpicture}
\end{document}
